% Implementation-reference-only options; no mathematical-guide dependency.
\section{Native implementation options}\label{sec:native-options}
\begingroup\small
\subsection{Build and select the full-index backend}
As of 19 October 2026, \code{calculate_index} and
\code{calculate_index_internal} default to \code{backend="native"}.
Native C++ performs expansion and singlet projection, with native LiE on
character-cache misses. The exact Sage Laurent-polynomial result, input format,
sector splitting and optional stored-sector reuse retain their existing contracts.
The implementation changes execution and storage, not the index algorithm.

\subsubsection{Build once for the Python interpreter used by Sage}
Run from the project root. The explicit build produces
\code{index._index_calculation_cpp}, including expansion, projection,
character caching and the LiE wrapper:
\begin{lstlisting}[language=bash]
sage -python -B index/cpp/build_index_calculation.py \
  --install-deps --cxx /usr/bin/g++
\end{lstlisting}
\code{--install-deps} installs pinned pybind11 and nlohmann/json headers in
\code{index/cpp/_build/deps}; omit it on subsequent builds. The build needs a
C++17 compiler, GMP/GMPXX and SQLite development files, plus Make and Bison
for the bundled LiE library. \code{--gmp-prefix} and \code{--sqlite-prefix}
select nonstandard dependency locations. Imports never compile or download.
Missing or outdated native builds report the build command; selecting the
native full-index backend does not silently fall back to FORM.

\subsubsection{Python API options}
\begin{tabularx}{\textwidth}{@{}p{49mm}Y@{}}
\toprule Option & Meaning for full-index calculation\\\midrule
\code{backend="native"} & Default: C++ expansion and projection; shared
character cache; no FORM expansion-cache reads or writes.\\
\code{backend="python"} & Original FORM/TFORM expansion and Python projection.
Its character-cache backend is selected independently (next page).\\
\code{worker_count=N} & Positive expansion thread count. Overrides
\code{form_threads}; if omitted, the existing \code{form_threads=1} default
applies to either backend.\\
\code{processes=N} & Python-backend character-cache process count. It does not
parallelize native LiE projection, which remains serial within each calculation.\\
\code{timeout=600} & Expansion and individual LiE-operation timeout in seconds.
\code{None} disables deadlines. Native checks are cooperative, with low-overhead
polling; this is not a deadline for the complete database/GUI job.\\
\code{char_cache_database_path} & Shared SQLite character/singlet cache file;
the default is \code{char_decomposition_cache.db} at the project root.\\
\code{form_cache_database_path} & Used only by the Python/FORM backend; defaults
beside the character cache. Native expansion is recomputed.\\\bottomrule
\end{tabularx}

For an existing valid theory dictionary \code{data}, the return type and
serialization stay the same:
\begin{lstlisting}[language=Python]
from index.n2_theory_index import calculate_index

value = calculate_index(data, 18, backend="native",
    worker_count=8, timeout=600,
    char_cache_database_path="characters.db")
reference = calculate_index(data, 18, backend="python",
    worker_count=8, processes=1,
    char_cache_database_path="characters.db")
\end{lstlisting}

\clearpage
\subsection{LiE, compatibility controls and command-line use}
The native backend calls \code{common/cpp/lie.cpp} through the bundled shared
library; no LiE executable is needed for the default
\code{lie_executable="lie"} setting. A custom LiE executable path selects
subprocess mode, preserving executable wrappers.

\subsubsection{Independent character-cache and LiE selections}
Set \code{N2SCFTDB_CHAR_CACHE_BACKEND} before importing the Python character
cache module:
\begin{tabularx}{\textwidth}{@{}p{28mm}Y@{}}
\toprule Value & Python character-cache behavior\\\midrule
\code{auto} & Default: use the C++ extension when present; use Python if it is
absent. Broken binaries or missing dependencies retain their import errors.\\
\code{cpp} & Require the C++ extension.\\
\code{python} & Use the original Python cache and executable LiE interface.\\\bottomrule
\end{tabularx}
This environment variable does not override \code{backend="native"} for full
indices. Selecting \code{backend="python"} alone does not force a Python cache.

The shared loader is \code{index._index_calculation_native.load_native}.
The full-index path passes
\code{require_index=True} to require expansion/projection support.

For \code{CharacterDecompositionCache} and \code{build_decomposition_cache},
the separate \code{lie_backend} option accepts \code{"auto"},
\code{"native"} or \code{"subprocess"}. Auto uses native LiE with the C++
cache and the default executable setting; a custom executable selects a
subprocess. Explicit native mode requires the extension and ignores the
executable setting. These constructor/builder options are not additional
\code{calculate_index} arguments.

\subsubsection{Command-line examples}
\begin{lstlisting}[language=bash]
# Native full-index calculation (also the default backend).
sage -python -B -m index.n2_theory_index theory.json \
  --order 18 --backend native --workers 8 \
  --char-cache-database characters.db

# Original Python cache, executable LiE and FORM/TFORM.
N2SCFTDB_CHAR_CACHE_BACKEND=python \
  sage -python -B -m index.n2_theory_index theory.json \
  --order 18 --backend python --workers 8 --processes 1
\end{lstlisting}
\code{--workers} overrides the legacy \code{--form-threads} option. For the
Python backend, one expansion thread selects \code{--form-executable}; larger
counts select \code{--tform-executable} with \code{-wN}. The native backend
does not use those executables or \code{--form-cache-database}.

\subsection{GUI behavior and engine logs}
The GUI uses native full-index calculation by default. Its existing
\code{tools/form_threads} preference also controls native expansion threads.
The CPU budget is divided by this thread count, rounding down, to allocate
at least one theory worker and at most one worker per selected job.
Native LiE projection is serial inside each theory worker. Coulomb-index
PE/PL still uses a separate FORM/TFORM executable; Coulomb spectrum calculation
uses neither FORM nor LiE.

Startup messages in the Index tab and \code{logs/log_index_*.log} identify
native expansion, shared-library or executable LiE, and separate FORM/TFORM
Coulomb expansion with its path and thread count. They describe selected
engines; cache or stored-sector reuse can avoid work. Spectrum-only runs
report that neither engine is used.
\endgroup
